%! Mode:: "TeX:UTF-8"
%! TEX program = xelatex
\PassOptionsToPackage{quiet}{xeCJK}
\documentclass[withoutpreface,bwprint]{cumcmthesis}
\usepackage{etoolbox}
\usepackage{ctex}
\BeforeBeginEnvironment{tabular}{\zihao{-5}}
\usepackage[numbers,sort&compress]{natbib}  % 文献管理宏包
\usepackage[framemethod=TikZ]{mdframed}  % 框架宏包
\usepackage{url}  % 网页链接宏包
\usepackage{subcaption}  % 子图宏包
\usepackage{tikz}
\usetikzlibrary{shapes.geometric, arrows.meta, positioning}

\newcolumntype{C}{>{\centering\arraybackslash}X}
\newcolumntype{R}{>{\raggedleft\arraybackslash}X}
\newcolumntype{L}{>{\raggedright\arraybackslash}X}
\renewcommand{\textbf}[1]{{\song\bfseries #1}}

\tikzstyle{startstop} = [rectangle, rounded corners, minimum width=3cm, minimum height=1cm,text centered, draw=black, fill=red!20]
\tikzstyle{process} = [rectangle, minimum width=3.5cm, minimum height=1cm, text centered, draw=black, fill=blue!20]
\tikzstyle{decision} = [diamond, aspect=2, minimum width=3.5cm, minimum height=1cm, text centered, draw=black, fill=green!20]
\tikzstyle{arrow} = [thick,->,>=stealth]

% ==================== 论文标题 ====================
\title{基于【核心模型或方法】的【研究对象】问题研究}

% ==================== 正文开始 ====================
\begin{document}

% 封面和目录页不显示页码
\pagenumbering{gobble}
\maketitle

% ==================== 摘  要 - BZD数模社制作,本文档为论文模板-如需详细说明可见论文模板-完整版====================
% ✓ 自查清单：
%   □ 第一段是否只用一句背景和一句问题引入，未堆砌技术细节？
%   □ 是否覆盖题干中的全部问题？
%   □ 是否根据真实关系衔接问题二、问题三和问题四，未强行写成递进模型？
%   □ 是否写出了能直接回答题目的量化结果？
%   □ 是否突出论文真实的创新点？
%   □ 是否使用第三人称并避免公式、图表和引用编号？
\begin{abstract}

	\noindent\textbf{摘要结构与内容要求：}

	% 第一段：引言段（背景+问题）
	随着【背景情境】的发展，【研究对象】对【现实目标或应用场景】具有重要影响。
	针对【研究对象】中存在的【核心问题】，本文围绕题目提出的若干任务展开研究。

	% 问题一：为解决【问题一的简化表述】，通过【关键分析或数据处理】确定【关键因素】，
	% 建立【模型名称】，将问题转化为【数学问题类型】，并采用【算法或软件】求解。
	% 得到【关键量化结果】，表明【明确结论】；通过【检验方法】验证，模型的【指标】为【具体数值】。

	针对问题一，为解决【问题一的简化表述】，通过【关键分析或数据处理】确定【关键因素】，
	建立【模型名称】，将问题转化为【数学问题类型】，并采用【算法或软件】求解。
	得到【关键量化结果】，表明【明确结论】；通过【检验方法】验证，模型的【指标】为【具体数值】。

	% 问题二：根据其与问题一的实际关系，【在问题一结果的基础上／结合前述指标／独立分析】
	针对问题二，为解决【问题二的简化表述】，根据其与问题一的实际关系，
	【在问题一结果的基础上／结合前述指标／独立分析】开展【关键处理】，
	建立【模型名称】，采用【算法】得到【核心结果】。
	结果显示【结论】，与【基准方案或实际数据】相比，【指标变化或误差】为【具体数值】。

	% 问题三
	针对问题三，为解决【问题三的简化表述】，结合【可用的前述结果或新增条件；若无直接联系则删除】，
	通过【关键分析】建立【模型名称】，采用【算法或软件】求解，得到【关键结果】。
	经【检验、对比或灵敏度分析】，说明【结果的可靠性及明确结论】。

	% 问题四（若题目有第四问）
	针对问题四，为解决【问题四的简化表述】，依据【题干新增条件或前述可用结果】，
	建立【模型名称】，将问题归结为【数学问题类型】，采用【算法】得到
	【方案、预测值、排序或决策结果】，最终提出【明确结论或建议】。

	% 最后一段：创新与评价
	本文模型在【创新点】方面具有一定特色。经【模型检验方式】验证，
	模型具有较好的【准确性、稳定性或鲁棒性】，可推广至【相似问题或应用领域】；
	针对【局限】，后续可通过【改进方向】进一步完善。

	\vspace{0.5cm}
	\keywords{【研究问题】；【主要模型】；【核心算法或创新方法】；【必要的补充关键词】}

\end{abstract}

% ==================== 问题重述 自查清单- BZD数模社制作,本文档为论文模板-如需详细说明可见论文模板-完整版 ====================
% ✓ 自查清单：
%   □ 是否完整说明了研究对象和问题背景？
%   □ 是否覆盖题目提出的全部问题？
%   □ 是否保留关键数据、时间范围、单位和限制条件？
%   □ 是否真正重组语言，而非仅替换同义词？
%   □ 是否删除了原题中的非必要图表和附件内容？
%   □ 是否没有提前写入模型、算法、求解过程和结果？
%   □ 研究背景是否简洁，未喧宾夺主？
%   □ 问题回顾是否使用明确的任务动词？
%   □ 所有参考文献和研究结论是否真实可核验？
%   □ 全章是否控制在1—2页或赛事规定范围内？

%%%%%%%%%%%%%%%%%%%%%%%%%%%%%%%%%%%%%%%%%%%%%%%%%%%%%%%%%%%%%
\pagenumbering{arabic}
\vspace{-1.4cm}

% ==================== 第一章：问题重述 ====================
\section{问题重述}

\subsection{研究背景}

【研究对象】是指【简要定义】，在【应用领域】中具有【主要作用】。
随着【社会、经济、技术或政策背景】的发展，【研究对象】呈现出【主要趋势】。
与此同时，其发展仍受到【主要矛盾或现实问题】的影响，因此有必要对【核心研究问题】进行系统研究。

\subsection{问题回顾}

根据题目提供的数据、条件和任务要求，需要解决以下问题：

\begin{enumerate}
	\item 分析【问题一研究对象】，确定【需要识别、评价或计算的内容】。
	\item 根据【数据条件】，建立模型完成【问题二的预测、评价或优化任务】。
	\item 考虑【新增条件】，研究【问题三的主要任务】。
	\item 结合【指定情境或可用的前述结果】，完成【问题四的任务】。
\end{enumerate}

\subsection{研究综述}

现有研究主要从【方向一】、【方向二】和【方向三】分析该问题。
在【方向一】方面，常用【模型或方法】研究【具体任务】；
在【方向二】方面，主要采用【模型或方法】分析【具体问题】；
在【方向三】方面，相关研究通常通过【模型或方法】评价【研究对象】。

现有方法为研究该问题提供了基础，但在【现有不足】方面仍具有一定局限。
基于题目条件和研究任务，有必要进一步从【本文建模角度】开展分析。

\newpage
%%%%%%%%%%%%%%%%%%%%%%%%%%%%%%%%%%%%%%%%%%%%%%%%%%%%%%%%%%%%%
% ==================== 第二章：问题分析- BZD数模社制作,本文档为论文模板-如需详细说明可见论文模板-完整版 ====================
% ✓ 自查清单：
%   □ 是否说明了题目整体需要完成什么？
%   □ 是否分析了各任务的数学类型？
%   □ 是否说明了各问题之间的数据或逻辑联系？
%   □ 是否分析了附件和补充数据的特点？
%   □ 是否说明了数据处理方法及选择理由？
%   □ 是否解释了模型选择依据和备选方案？
%   □ 是否说明了模型检验思路？

\section{问题分析}

题目围绕【研究对象】提出了【概括全部任务】。从数学建模角度看，各问题分别涉及
【评价、预测、分类、优化或决策等类型】，并通过【数据、指标、参数或中间结果】形成较强的逻辑联系。
前序问题主要用于完成【基础任务】，所得【结果类型】将为后续【任务】提供输入；
后续问题则进一步考虑【新增条件】，最终形成由【起点任务】到【终点任务】的完整分析过程。

题目提供的数据主要包括【数据内容】，同时需要补充收集【外部数据】。
考虑到数据可能存在【缺失、异常、量纲不同、统计口径不一致或指标相关性较强】等问题，
需要先通过【数据处理方法】完成数据清洗与口径统一，再采用
【标准化、相关性分析、主成分分析或特征筛选方法】提取有效信息，为后续建模提供可靠数据基础。

对于前序问题，需要解决【任务概括】，其本质属于【数学问题类型】。
根据【数据特点和题目条件】，可采用【模型或方法】刻画【变量之间的关系】，
并将得到的【权重、参数、关键指标或分类结果】传递至后续问题。
为降低单一模型的局限，可以选择【备选方法】，并依据【评价标准】进行比较。

在前述分析的基础上，后续问题需要完成【预测、优化、影响评价或情景分析】。
针对【预测任务】，可利用【候选方法】描述变化规律，并根据【误差指标】选择模型或进行组合；
针对【优化或决策任务】，可将前述结果作为模型参数，以【目标】构建【规划或决策模型】，
并设置【主要约束】；对于【外部政策、环境变化或不确定性】，
可通过参数调整和情景设置分析其影响。

为检验模型的合理性，可采用【误差分析、交叉验证、对比实验、灵敏度分析或稳健性检验】
评价模型表现。最终输出【题目要求的预测值、评价结果、优化方案或建议】。
% ==================== 整体求解框架图 ====================
\begin{figure}[htb]
	\centering
	\includegraphics[width=\linewidth]{5e134de4bc90d519e67cad3bb4f51c93.png}
	\caption{流程框架图}
\end{figure}

\newpage
%%%%%%%%%%%%%%%%%%%%%%%%%%%%%%%%%%%%%%%%%%%%%%%%%%%%%%%%%%%%%
% ==================== 第三章：模型假设- BZD数模社制作,本文档为论文模板-如需详细说明可见论文模板-完整版====================
% ✓ 自查清单：
%   □ 是否包含题目明确给出的关键条件？
%   □ 是否包含后续模型真正需要的假设？
%   □ 每条假设是否有合理依据？
%   □ 是否说明主要简化和研究边界？
%   □ 是否保留核心因素，只忽略次要因素？
%   □ 数据假设是否与实际数据处理一致？
%   □ 是否存在重复或完全没有在后文使用的假设？

\section{模型假设}

为简化问题，本文做出以下假设：

\begin{enumerate}[itemindent=2em]
	\item 假设【具体内容】。

	【说明假设依据、简化对象及其对模型的作用。】

	\item 假设【具体内容】。

	【说明假设依据、简化对象及其对模型的作用。】

	\item 假设【具体内容】。

	【说明假设依据、简化对象及其对模型的作用。】

	\item 假设【具体内容】。

	【说明假设依据、简化对象及其对模型的作用。】

	\item 假设【具体内容】。

	【说明假设依据、简化对象及其对模型的作用。】
\end{enumerate}

\newpage
%%%%%%%%%%%%%%%%%%%%%%%%%%%%%%%%%%%%%%%%%%%%%%%%%%%%%%%%%%%%%
% ==================== 第四章：符号说明 ====================
% ✓ 自查清单：
%   □ 是否采用三线表？
%   □ 是否只收录反复使用的主要符号？
%   □ 是否删除只出现一次的临时变量？
%   □ 有量纲变量是否标明单位？
%   □ 同一符号是否只表示一个含义？
%   □ 同一变量是否始终使用同一符号和单位？
%   □ 是否检查了符号冲突和单位不一致？

\section{符号说明}

为便于模型表达，现将全文反复使用的主要符号汇总如表 \ref{tab:符号说明} 所示。

\begin{table}[H]
	\centering
	\caption{主要符号说明}
	\begin{tabularx}{\textwidth}{C|L|C}
		\toprule
		\textbf{符号} & \textbf{含义} & \textbf{单位} \\
		\midrule
		$n$ & 样本总数 & 个 \\
		$d$ & 每个样本的特征变量个数 & — \\
		$\boldsymbol{x}^{(i)}$ & 第 $i$ 个样本的 $d$ 个特征变量，$\boldsymbol{x}^{(i)} = (x_1^{(i)}, x_2^{(i)}, \dots, x_d^{(i)})$ & — \\
		$y^{(i)}$ & 第 $i$ 个样本对应的目标变量 & — \\
		$f(\cdot)$ & 构建的预测函数 & — \\
		$\Omega(f)$ & 正则化项 & — \\
		【符号】 & 【含义】 & 【单位】 \\
		\bottomrule
	\end{tabularx}
	\label{tab:符号说明}
\end{table}

\noindent\textbf{注：}表中未列出的局部符号将在正文首次出现时说明。

\newpage
%%%%%%%%%%%%%%%%%%%%%%%%%%%%%%%%%%%%%%%%%%%%%%%%%%%%%%%%%%%%%
% ==================== 第五章：问题一的模型建立和求解 ====================
% ✓ 自查清单：
%   □ 是否针对当前题目重新定义变量、目标函数、约束和参数？
%   □ 前一问题输出是否真实进入后一问题？符号、单位是否一致？
%   □ 参数是否有来源？公式是否有引出和解释？
%   □ 求解前是否集中给出目标函数、约束条件和变量范围？
%   □ 是否说明算法思想、参数、初值、迭代次数？
%   □ 是否给出具体数值并逐项回答题目？

\section{问题一的模型建立和求解}

\subsection{建模思路与模型选择}

问题一要求【问题目标】。根据题目条件和数据特征，该问题在数学上属于【问题类型】。
考虑到【模型适用理由】，选择【模型名称】进行求解，并针对【本题特点】对模型的
【目标函数、约束、参数或算法】进行调整。

\subsection{数据处理}

本问题使用【数据来源】。首先进行【处理方法】，原因是【处理理由】。
处理后保留【样本数量】个有效样本，并构造【特征名称】作为模型输入。

\subsection{模型的建立}

定义【核心变量】和【关键参数】。根据【理论依据或现实机理】建立【目标函数或核心方程】，
并结合【资源、时间、容量等条件】设置约束。

综合上述分析，建立如下模型：

\begin{equation}
	\text{minimize} \quad F(x) = \text{【目标函数】}
\end{equation}

\begin{equation}
	\text{subject to} \quad \left\{
	\begin{array}{l}
		\text{【约束条件1】} \\
		\text{【约束条件2】} \\
		\vdots \\
		\text{【约束条件n】}
	\end{array}
	\right.
\end{equation}

\subsection{模型的求解}

采用【算法名称】求解。算法首先……，其次……，最后……。
使用【软件与版本】，关键参数为……，最大迭代次数为……，终止条件为……。
完整程序见附件【编号】。

\subsection{结果分析}

求解得到【核心结果】。由图 \ref{fig:xxx} 或表 \ref{tab:xxx} 可知……。
其中，【关键变量】对结果影响最大，【主要趋势】与实际规律一致，
最终得到问题一的答案为【具体结论】。

\subsection{模型检验与灵敏度分析}

采用【检验方法】验证模型，指标【名称】为……，满足【标准】。
进一步改变参数【名称】，结果表明模型在【区间】内保持稳定，主要结论未发生改变。

\newpage
%%%%%%%%%%%%%%%%%%%%%%%%%%%%%%%%%%%%%%%%%%%%%%%%%%%%%%%%%%%%%
% ==================== 第六章：问题二的模型建立和求解 ====================
\section{问题二的模型建立和求解}

\subsection{建模思路与模型选择}
【内容结构参考问题一】

\subsection{数据处理}
【内容结构参考问题一】

\subsection{模型的建立}
【内容结构参考问题一】

\subsection{模型的求解}
【内容结构参考问题一】

\subsection{结果分析}
【内容结构参考问题一】

\subsection{模型检验与灵敏度分析}
【内容结构参考问题一】

\newpage
%%%%%%%%%%%%%%%%%%%%%%%%%%%%%%%%%%%%%%%%%%%%%%%%%%%%%%%%%%%%%
% ==================== 第七章：问题三的模型建立和求解 ====================
\section{问题三的模型建立和求解}

\subsection{建模思路与模型选择}
【内容结构参考问题一】

\subsection{数据处理}
【内容结构参考问题一】

\subsection{模型的建立}
【内容结构参考问题一】

\subsection{模型的求解}
【内容结构参考问题一】

\subsection{结果分析}
【内容结构参考问题一】

\subsection{模型检验与灵敏度分析}
【内容结构参考问题一】

\newpage
%%%%%%%%%%%%%%%%%%%%%%%%%%%%%%%%%%%%%%%%%%%%%%%%%%%%%%%%%%%%%
% ==================== 第八章：问题四的模型建立和求解 ====================
\section{问题四的模型建立和求解}

\subsection{建模思路与模型选择}
【内容结构参考问题一】

\subsection{数据处理}
【内容结构参考问题一】

\subsection{模型的建立}
【内容结构参考问题一】

\subsection{模型的求解}
【内容结构参考问题一】

\subsection{结果分析}
【内容结构参考问题一】

\subsection{模型检验与灵敏度分析}
【内容结构参考问题一】

\newpage
%%%%%%%%%%%%%%%%%%%%%%%%%%%%%%%%%%%%%%%%%%%%%%%%%%%%%%%%%%%%%
% ==================== 模型分析与检验 - BZD数模社制作,本文档为论文模板-如需详细说明可见论文模板-完整版====================
% ✓ 自查清单：
%   □ 是否使用匹配的量化指标，并给出判定标准？
%   □ 灵敏度分析的参数选择和变化范围是否有依据？
%   □ 是否分析输入扰动、假设变化、误差来源和传播？

\section{模型的分析与检验}

\subsection{灵敏度分析}

为了评估模型对关键参数变化的响应程度，对参数【参数名】进行灵敏度分析。

改变参数【参数名】在【变化范围】内变化，观察目标函数【目标函数】的变化情况。
结果表明……（描述敏感性结论）。

\subsection{误差分析}

模型的预测误差主要来自【误差来源】。
采用【误差衡量方法】，计算得误差为【误差数值】，模型的【评价指标】满足要求。

\newpage
%%%%%%%%%%%%%%%%%%%%%%%%%%%%%%%%%%%%%%%%%%%%%%%%%%%%%%%%%%%%%
% ==================== 模型评价、改进与推广- BZD数模社制作,本文档为论文模板-如需详细说明可见论文模板-完整版 ====================
% ✓ 自查清单：
%   □ 是否根据前文编号设置为第九章或第十章？
%   □ 是否用一段话概括主要模型、联动关系和核心结论？
%   □ 每项优点是否有具体做法或量化证据？
%   □ 缺点是否说明来源、影响和适用范围？
%   □ 改进方向是否逐项对应模型缺点？
%   □ 是否区分模型改进与模型推广？

\section{模型的评价、改进与推广}

本文针对【研究问题】，分别建立【模型一】、【模型二】和【模型三】。
各模型依次完成【问题一目标】、【问题二目标】和【问题三目标】，
并通过【检验方法】验证结果可靠性。现对模型的优点、不足、改进方向与推广价值进行总结。

\subsection{模型的优点}

\begin{enumerate}[itemindent=2em]
	\item 模型具有较强针对性。针对【题目特点】，在【经典模型】基础上增加【变量或约束】，
	提高了模型与实际问题的匹配程度。

	\item 数据处理较为完整。针对【数据问题】采用【处理方法】，
	降低了数据质量对后续模型的影响。

	\item 问题之间联动清晰。问题一得到的【结果】作为问题二的【输入】，
	并进一步进入问题三的【目标函数或约束】。

	\item 模型结果可靠。经【检验方法】验证，【量化指标】达到【具体数值】，
	主要结论与实际规律一致。
\end{enumerate}

\subsection{模型的不足}

\begin{enumerate}[itemindent=2em]
	\item 受【数据来源】限制，未能充分考虑【因素】，可能对【具体结果】产生一定影响。

	\item 为保证模型可解，对【现实过程】进行了【具体简化】，
	使模型在【情景】下的适用性受到限制。

	\item 模型结果对参数【名称】较为敏感，该参数估计误差可能影响最终结论。
\end{enumerate}

\subsection{模型的改进与推广}

\begin{enumerate}[itemindent=2em]
	\item 进一步收集【数据类型】，扩大时间和空间范围，提高参数估计可靠性。

	\item 增加【变量或约束】，进一步描述当前未充分考虑的现实因素。

	\item 引入【动态模型、鲁棒优化或改进算法】，降低不确定性对结果的影响。

	\item 模型可推广到【对象或场景】；推广时需要重新标定【参数】，替换【变量】，并增加【新约束】。
\end{enumerate}

\newpage
%%%%%%%%%%%%%%%%%%%%%%%%%%%%%%%%%%%%%%%%%%%%%%%%%%%%%%%%%%%%%
% ==================== 参考文献 ====================

\section*{AI工具使用声明}

\noindent 【选择以下之一：】

\noindent \textbf{情形A：未使用AI工具}

本参赛队在竞赛过程中未使用任何AI工具。

\vspace{0.5cm}
\noindent \textbf{情形B：使用了AI工具}

本参赛队在竞赛过程中使用了AI工具，主要用于【文献检索辅助、建模思路讨论、数据处理建议、程序代码调试、结果核验、图表优化、语言润色和格式检查等实际用途】，详细使用情况见支撑材料《AI工具使用详情.pdf》。

\newpage

\noindent 参考文献应按照在正文中首次出现的顺序编号。

\begin{thebibliography}{99}
	\bibitem{ref1} 作者名. 文章标题[J]. 期刊名称, 年份, 卷(期): 起止页码.

	\bibitem{ref2} 作者名. 书名[M]. 出版地: 出版社, 出版年.

	\bibitem{ref3} 作者名. 学位论文标题[D]. 授予单位, 年份.

	% 禁止事项：不能将 ChatGPT、DeepSeek、Claude、GitHub Copilot 等 AI 工具列为参考文献
\end{thebibliography}

\newpage
%%%%%%%%%%%%%%%%%%%%%%%%%%%%%%%%%%%%%%%%%%%%%%%%%%%%%%%%%%%%%
% ==================== AI工具使用声明 ====================

%%%%%%%%%%%%%%%%%%%%%%%%%%%%%%%%%%%%%%%%%%%%%%%%%%%%%%%%%%%%%
% ==================== 附  录 ====================

\appendix

\section{支撑材料文件列表}

\begin{table}[H]
	\centering
	\begin{tabularx}{\textwidth}{L|L}
		\toprule
		\textbf{文件名} & \textbf{功能描述} \\
		\midrule
		% 数据处理
		数据预处理.py & 数据导入、清洗、缺失值处理和标准化 \\
		数据分析.ipynb & 数据统计特征分析和可视化 \\

		% 问题一程序
		problem1\_model.py & 问题一模型建立与求解 \\
		problem1\_validation.py & 问题一模型检验与灵敏度分析 \\

		% 问题二程序
		problem2\_model.py & 问题二模型建立与求解 \\

		% 问题三程序
		problem3\_model.py & 问题三模型建立与求解 \\

		% 问题四程序
		problem4\_model.py & 问题四模型建立与求解 \\

		% 辅助材料
		AI工具使用详情.pdf & AI工具在竞赛中的具体使用情况说明 \\

		\bottomrule
	\end{tabularx}
	\label{tab:文件列表}
\end{table}

\section{计算环境与求解工具}

\noindent 本文所用的编程环境及相关工具如下：

\begin{itemize}
	\item \textbf{编程语言：}Python 3.9+
	\item \textbf{主要库与工具：}
	\begin{itemize}
		\item 数据处理：Pandas, NumPy
		\item 科学计算：SciPy
		\item 机器学习：Scikit-learn
		\item 可视化：Matplotlib, Seaborn
		\item 优化求解：CVXPY, PuLP
	\end{itemize}
	\item \textbf{运行环境：}Windows 10/11 或 Linux, 内存 8GB 以上
\end{itemize}

\section{源程序代码}

\noindent 各问题的完整源程序如下（详细注释见相应.py文件）：

% ========== 问题一源代码 ==========
\subsection{问题一源程序}

\noindent \textbf{程序名称：}problem1\_model.py

\noindent \textbf{程序功能：}完成问题一的数据输入、模型建立和求解。

\noindent \textbf{主要输入：}【输入数据文件名、格式和说明】

\noindent \textbf{主要输出：}【输出结果文件名和格式】

\noindent \textbf{运行方式：}
\begin{verbatim}
python problem1_model.py
\end{verbatim}

% 代码内容（可选，建议放在单独的代码文件中，这里仅示意框架）
\noindent \textbf{核心代码框架：}
\begin{verbatim}
# 导入必要的库
import pandas as pd
import numpy as np
from sklearn import ...

# 加载数据
data = pd.read_csv('data.csv')

# 数据处理
# ...

# 模型建立和求解
# ...

# 结果输出
# ...
\end{verbatim}

% ========== 问题二源代码 ==========
\subsection{问题二源程序}

\noindent \textbf{程序名称：}problem2\_model.py

\noindent 【按照问题一的格式填写】

% ========== 问题三源代码 ==========
\subsection{问题三源程序}

\noindent \textbf{程序名称：}problem3\_model.py

\noindent 【按照问题一的格式填写】

% ========== 问题四源代码 ==========
\subsection{问题四源程序}

\noindent \textbf{程序名称：}problem4\_model.py

\noindent 【按照问题一的格式填写】

\section{补充数据与计算结果}

\noindent 本章节列出正文未完整展示的长表、推导、参数和补充图表。

\begin{table}[H]
	\centering
	\caption{补充数据表}
	\begin{tabularx}{\textwidth}{C|C|C|C}
		\toprule
		\textbf{变量} & \textbf{数值1} & \textbf{数值2} & \textbf{数值3} \\
		\midrule
		【行标签】 & 【数据】 & 【数据】 & 【数据】 \\
		\bottomrule
	\end{tabularx}
	\label{tab:补充数据}
\end{table}

\end{document}
